\input{preamble.tex}
\title{Concrete markings in the amalgamation examples}
\author{Human Compression project}
\date{Private candidate, 11 October 2026}
\begin{document}\maketitle
This note identifies the remaining concrete markings in
Piedade--Tranchida, arXiv:2505.24662v1. It uses the reviewed tree,
indexed HNN and orbit-transfer theorems with their exact hypotheses.
The retained normal-form input is R.~C.~Lyndon--P.~E.~Schupp,
\emph{Combinatorial Group Theory} (2001 edition), Chapter I,
Theorem 1.2 for free products, and Chapter IV, Theorem 2.6 and
Britton's lemma for amalgams and HNN extensions, as cited by the
source. Bass--Serre tree realizations and the Coxeter reflection
models are classical inputs, not new conclusions of finite examples.

\section{The small free products}
Types remain part of a coset system even if they contain one element.
Thus Example 3.10's rank-one factor $(1,(1))$ produces the original
parabolics together with one equal to the whole group. The result
is the cone on the input, with one new universal element, and is
not type-preservingly isomorphic to the input. A rank-zero factor
would give the asserted identity. In Example 3.11 the two factors
$(C_2,(1))$ give the Bass--Serre tree of $C_2*C_2$: its vertices
are cosets of the two factors, its edges are group elements, and
both vertex degrees are two. This is the infinite bivalent tree.

For an $n$-gon ($n\geq3$) and the rank-one $C_2$ factor, the
presentation is
$\langle a_1,a_2,a_3\mid a_i^2=1,(a_1a_2)^n=1\rangle$.
The subgroups obtained by deleting one generator are precisely
Example 3.12's maximal parabolics. Under the classical Coxeter
reflection-triangle realization it is the tessellation with angles
$\pi/n,0,0$. Amalgamating an $n$-gon and an $m$-gon over their
shared reflection instead gives
\[
 \langle a_1,a_2,b_3\mid a_1^2=a_2^2=b_3^2=1,
 (a_1a_2)^n=(a_2b_3)^m=1\rangle.
\]
Its maximal parabolics are the three deleted-generator subgroups,
as in Example 3.26. For $n,m\geq3$ the angles $\pi/n,\pi/m,0$
define a hyperbolic triangle. More generally this interpretation
requires $1/n+1/m<1$; the algebraic amalgam formula has no such
geometric inequality.

Example 3.25's empty shared type set means amalgamation over the
identified Borel subgroups, and is a free product if both are
trivial. Identifying the two entire systems gives that system.
The parenthetical assertion that thin flag-transitive systems always
have trivial Borel requires a faithful acting group. It is false
for arbitrary specified group actions; the reviewed exact panel
formula keeps the Borel explicitly.

\section{The amalgamation map is part of compatibility}
The positive amalgamation theorem uses one chosen isomorphism of the
common parabolics carrying every indexed intersection to the matching
one. Separate abstract isomorphisms, as written in Definition 3.16,
do not imply this coherence for an arbitrary chosen amalgamation.
For example take two standard rank-two geometries of $C_2^2$, with
coordinate generators $a,b$ and corresponding maximal subgroups
$\langle a\rangle,\langle b\rangle$. Let all types be shared, so
the amalgamated subgroup is the whole group, but identify the two
copies by the automorphism interchanging $a,b$. Every pair appearing
in Definition 3.16 is abstractly isomorphic. Nevertheless the output
maximal subgroup of either type is
$\langle a,b\rangle=C_2^2$, since it combines different coordinate
lines under the chosen identification. The output consists of two
universal singleton types and is not firm, while both inputs are thin.

Thus Lemma 3.17's step from abstract isomorphism to equality of the
embedded intersections needs the specified coherent map. The explicit
assumption on page 23 that all the maximal parabolics are special
supplies exactly that condition. All positive uses here retain it;
none applies a gluing theorem to arbitrary abstractly isomorphic
parabolic data. This example qualifies the literal shorthand, without
refuting the coherent-map theorem proved in the common tree note.

\section{The incompatible $S_4$ parabolics}
In Example 3.18 the literal generators $a_0,a_1,a_2$ satisfy the
Coxeter $[4,3]$ relations and $(a_0a_1a_2)^3=1$. In the signed
permutation model of $[4,3]$, of order 48, this cube of a Coxeter
element is central inversion. Hence $A_3=\langle a_0,a_1,a_2\rangle$
has order at most 24. The assignment
\[
 a_0=(12)(34),\quad a_1=(23),\quad a_2=(34),\quad a_3=(14)(23)
\]
satisfies every relator of the full source group in $S_4$. The first
three images generate $S_4$, proving $A_3\cong S_4$ exactly.
The pair $a_0,a_1$ generates $D_8$: its defining order bound is
attained in this quotient. Since $a_3$ commutes with both,
$A_2=\langle a_0,a_1,a_3\rangle$ has order dividing 16.
Lagrange's theorem now bounds $|A_2\cap A_3|$ by
$\gcd(16,24)=8$, and this intersection contains the displayed
$D_8$. Thus equality holds. On the other side the Coxeter $[3,3,4]$
parabolic $B_4$ is $[3,3]\cong S_4$ and
$B_2\cap B_4=\langle b_0,b_1\rangle\cong D_6$.
These different intersection orders prove the claimed incompatibility
without retaining the classification of the whole regular polytope.

\section{The simplex HNN construction}
Take $a_i=(i,i+1)$ in $S_6$, $1\leq i\leq5$, as in
Example 3.36. The associated subgroups are
$C=\langle a_2,a_3\rangle$ and $D=\langle a_3,a_4\rangle$,
with indexed map $a_2\mapsto a_3$, $a_3\mapsto a_4$.
The type classes are $\{1\},\{2,3,4\},\{5\}$ and $F=\{4\}$.
The reviewed indexed HNN theorem therefore gives
\[
 G_1=\langle A_1,t\rangle,\quad G_{234}=\langle a_1,a_5,t\rangle,
 \quad G_5=\langle A_5,t\rangle,\quad
 G_t=A_4=\langle a_1,a_2,a_3,a_5\rangle.
\]
Indexed compatibility follows from the named bijection and classical
Coxeter parabolic embeddings. The resulting geometry is
flag-transitive, residually connected and firm. Its Borel is trivial,
and the exact panel formula gives $2,6,2,\infty$.
Example 3.37's whole-base automorphism gives $A\rtimes\ZZ$ with
the stated orbit-coarsened parabolics when its indexed action is
specified. For firm inputs the reviewed Boolean-parabolic lemma
supplies the indexed action from source admissibility. Preservation
of an unlabelled parabolic family alone does not supply it on an
arbitrary nonfirm input.

\section{The tetrahedral Artin example}
In Example 5.7 let $A$ be the Artin group on $a_1,a_2,a_3,a_4$
with all pair labels three, and let $B=\langle p=(23),q=(34)\rangle$
act on the names. At representative 2 the orbits of
$1,\langle q\rangle,\langle p\rangle,B$ are
$\{2\},\{2\},\{2,3\},\{2,3,4\}$. These give the exact
orbit-intersection condition; the fixed singleton orbit $\{1\}$
also satisfies it. The complete orbit-transfer theorem applies.
The maximal parabolics are
\[
 A(\{2,3,4\})\rtimes B,\quad A(\{1\})\rtimes B,\quad
 A(\{1,2\})\rtimes\langle q\rangle,\quad
 A(\{1,2,3\})\rtimes\langle p\rangle.
\]
They are the deleted-generator subgroups of $(a_1,a_2,p,q)$.
This proves the source's firm, residually connected, flag-transitive
rank-four geometry; its panels have sizes $\infty,\infty,2,2$.

There is also a four-generator presentation, with $a=a_1,x=a_2$:
\[
 \begin{gathered}
 p^2=q^2=(pq)^3=1,\qquad [a,p]=[a,q]=[x,q]=1,\\
 \W_3(a,x)=\W_3(x,a),\qquad\W_6(x,p)=\W_6(p,x).
 \end{gathered}
\]
Set $x_2=x$, $x_3=pxp$, $x_4=qpxpq$. The relation $pqp=qpq$
and $[x,q]=1$ show that $p$ swaps $x_2,x_3$ and fixes $x_4$,
while $q$ swaps $x_3,x_4$ and fixes $x_2$.
The doubled-braid identity converts the length-six relation into
the length-three braid on $x_2,x_3$. Conjugation supplies the other
pair braids and the braids with $a$. These formulas give mutually
inverse maps with the original six-generator semidirect presentation.

\section{The three spanning trees in Example 5.18}
Use the six source generators and stable letters $t_G,t_H,t_K$.
For the first tree the relations are
\[
 a_1=b_4,\quad b_3=c_6,\quad t_G^{-1}c_5t_G=a_2,
 \quad t_G^{-1}c_6t_G=a_1.
\]
Eliminate four names, obtaining the free basis
$x=b_3=c_6$, $y=c_5$, $t=t_G$, with
$a_1=b_4=t^{-1}xt$ and $a_2=t^{-1}yt$.
For the second tree use $x=a_1=b_4=c_6$, $y=a_2=c_5$, $t=t_H$,
and eliminate $b_3=t^{-1}xt$. For the third use
$x=a_1=c_6=b_3$, $y=a_2=c_5$, $t=t_K$,
and eliminate $b_4=t^{-1}xt$.
Every substitution is reversible and leaves no relator on $x,y,t$;
these are explicit isomorphisms onto $F(x,y,t)$.

For each tree the final types are $X=\{1,3,4,6\}$,
$Y=\{2,5\}$ and the stable type. Their maximal parabolics become
\[
 \langle y,t\rangle,\qquad\langle x,t\rangle,\qquad\langle x,y\rangle.
\]
For instance the first tree's $X$ subgroup is
$\langle a_2,c_5,t_G\rangle=\langle t^{-1}yt,y,t\rangle
=\langle y,t\rangle$, and its stable subgroup is
$\langle b_3,c_5,c_6\rangle=\langle x,y\rangle$.
All others follow from the displayed substitutions. Thus the same
group isomorphisms identify the actual parabolics and induce coset
incidence isomorphisms, establishing the full example. Each is the
standard rank-three free-group geometry, with infinite panels.
On source page 56, $a_2=c_5$ is incorrectly claimed also for the
first tree; there they are conjugate. Merely finding separately
isomorphic parabolic groups would not suffice for an incidence
isomorphism, whereas the explicit marked maps do.

This example does not settle the general fundamental-geometry or
universal-property questions in Sections 5.4--6. The separate reviewed
obstruction concerns only the naive sequential algorithm. All source
application attributions remain with Piedade and Tranchida; no
historical-priority or new-question claim is made here.
\end{document}
